% SGA 3, 第 VI_B 篇, 第 1 轮英文重构底本的中译.
% 范围: 引理 8.1.2 及其完整证明.
% 权威底本: Polo--Gille Exp6B-13oct24.pdf, 组件第 61--62 页
% (合卷第 402--403 页; 重版页眉第 391--392 页; 原版边栏页码第 375--376 页).
% 命题 8.2 始于组件第 62 页, 不在本组件内.
% OCR 与 Reinhold 英文 Markdown 仅为对照材料.
% 第 2 轮已将原复合态射裁图替换为忠实的原生 LaTeX 重构, 替换已完成.

\medskip
\phantomsection\label{lem:sga3-VIB-8.1.2}%
\noindent\textbf{引理 8.1.2.---}
\emph{设 \(G\) 是一个群层, \(A,B\) 是 \(G\) 的两个群子预层.}
\begin{enumerate}
\item[\textup{\SGAInlineRefTarget{lem:sga3-VIB-8.1.2:item:i}(i)}]\phantomsection\label{concl:sga3-VIB-8.1.2-i}
  \emph{态射}
  \[
    \tau:\operatorname{Comm}(A,B)
      \longrightarrow
      \operatorname{Comm}(A^\flat,B^\flat)
  \]
  \emph{是一个覆盖单态射.}

\item[\textup{(ii)}]\phantomsection\label{concl:sga3-VIB-8.1.2-ii}
  \emph{因此有同构}
  \phantomsection\label{form:sga3-VIB-8.1.2-CommT-isomorphism}%
  \[
    \operatorname{Comm}_T(A,B)
      \xrightarrow{\ \sim\ }
    \operatorname{Comm}_T(A^\flat,B^\flat).
  \]
\end{enumerate}

\emph{证明.}
\textup{\SGARefLink{concl:sga3-VIB-8.1.2-i}{(i)}}
由于 \(A\) (分别地, \(B\)) 是 \(G\) 的子预层,
相伴层 \(A^\flat\) (分别地, \(B^\flat\))
是 \(G\) 的一个包含 \(A\) (分别地, \(B\)) 的子预层.
所以 \(\tau\) 是单态射.

现在证明 \(\tau\) 是覆盖.
设 \(S\in\operatorname{Ob}\mathcal C\), 且
\[
  g\in\operatorname{Comm}(A^\flat,B^\flat)(S).
\]
存在一个整数 \(n\geq1\), 以及对 \(1\leq i\leq n\), 元素
\[
  a'_i\in A^\flat(S),\qquad
  b'_i\in B^\flat(S),\qquad
  \varepsilon_i\in\{\pm1\},
\]
使得
\[
  g=(a'_1,b'_1)^{\varepsilon_1}\cdots
    (a'_n,b'_n)^{\varepsilon_n},
\]
其中 \((a,b)\) 表示换位子 \(aba^{-1}b^{-1}\).
存在 \(S\) 的一个加细 \(R\), 使得对每个 \(i\), 有
\[
  a'_i\in A(R)
  \qquad\text{及}\qquad
  b'_i\in B(R)
\]
于是 \(g\) 是复合态射
\begin{center}
  \phantomsection\label{diag:sga3-VIB-8.1.2-composite}%
  \begin{tikzcd}[column sep=large]
    R
      \arrow[r,"{(a'_1,\ldots,b'_n)}"] &
    (A\times B)^n
      \arrow[r,"{\Phi^{\varepsilon_1,\ldots,\varepsilon_n}}"] &
    \operatorname{Comm}(A,B)
  \end{tikzcd}
\end{center}
其中 \(\Phi=\Phi^{\varepsilon_1,\ldots,\varepsilon_n}\)
是在集合上对每个 \(T\in\operatorname{Ob}\mathcal C\),
\(a_i\in A(T)\) 与 \(b_i\in B(T)\), 由下式定义的态射:
\phantomsection\label{form:sga3-VIB-8.1.2-Phi}%
\[
  \Phi(a_1,b_1,\ldots,a_n,b_n)
    =(a_1,b_1)^{\varepsilon_1}\cdots
     (a_n,b_n)^{\varepsilon_n}
\]
这表明
\[
  g\in\operatorname{Comm}(A,B)(R).
\]
于是, 如 IV,
\SGARefLink{sga3:iv:proposition:4:3:11}{4.3.11}\textup{\SGARefLink{lem:sga3-VIB-8.1.2:item:i}{(i)}}
的证明那样, 可知
\[
  \operatorname{Comm}(A,B)
    \longrightarrow
  \operatorname{Comm}(A^\flat,B^\flat)
\]
是覆盖.

态射
\[
  \operatorname{Comm}(A^\flat,B^\flat)
    \longrightarrow
  \operatorname{Comm}(A^\flat,B^\flat)^\flat
\]
同样是覆盖单态射
(IV, \SGARefLink{sga3:iv:proposition:4:3:11}{4.3.11}\textup{\SGARefLink{sga3:iv:proposition:4:3:11:item:iv:01}{(iv)}}).
所以
\[
  \operatorname{Comm}(A,B)
    \longrightarrow
  \operatorname{Comm}(A^\flat,B^\flat)^\flat.
\]
也如此. 因而由 IV, \SGARefLink{sga3:iv:corollary:4:3:12}{4.3.12}, 得到同构
\[
  \operatorname{Comm}(A,B)^\flat
    \xrightarrow{\ \sim\ }
  \operatorname{Comm}(A^\flat,B^\flat)^\flat.
\]
\hyperref[lem:sga3-VIB-8.1.2]{引理 8.1.2} 得证.

% 中文翻译检查点：本组件全部正文、证明、脚注及源注已译；源119行。
% 机械核对：行内数学26处、陈列数学12处及全部引用键守恒（数学中的自然语言已译）。
% 状态：完成源文覆盖和逐段自校；尚未经独立合卷编译与页面审查。
